\scanpage{42}
\unsure{2)|2°)|$\beta$)} Cas $n$ impair $\rho=\widetilde\rho$, $\sigma=2\widetilde\sigma$, l'équation s'écrit

\[
\lambda\widetilde\sigma(2+\lambda\rho)=0
\]

Donc si $\sigma\inserted{\unsure{\sigma}}(=2\widetilde\sigma)$ et $\rho\inserted{(=\widetilde\rho)}$ inv., \deleted{la seule} $\exists\unsure{1|!}$ section $u$ de $U_\rho$, partout $\ne1$, \deleted{telle} qui invarie $f$, ou ce qui revient au $=$, qui invarie $Q$ -- savoir
\margin{\uline{NB} Si $\widetilde\sigma$ inv. et non $\sigma$ (p. \unsure{x.|ex.} en car. 2) on a \unsure{quand même} unique solution en $\lambda$, savoir $\lambda=-\dfrac{2}{\rho}$, mais celle-ci peut être nulle en certains pts. \unsure{Donner|Donnons|Donc} énoncé commun aux cas 1°) et \unsure{2°)} \unsure{en termes de}  $\boxed{\sigma\widetilde\rho\ \text{inversibles}}$}
% slanted

\[
\underbrace{u=\mathrm{id}\unsure{-|+}\frac{2}{\rho}\pi'\otimes\pi}_{\text{antipodisme}}
\]

D'autre part, conditions équiv.\deleted{\unsure{alentes}}

\deleted{\unsure{si}} a) $2\widetilde\sigma=\sigma=0$ et $\widetilde\sigma\rho\inserted{(=\unsure{\widetilde\sigma\widetilde\rho|\sigma\rho})}=0$.

[\uline{NB} si 2 inv., celà équivaut à $\sigma=0$; si $\unsure{2\cdot1_S|2\,1_S}=0$, celà équivaut à $\widetilde\sigma\rho=0$]

b) $U_\rho$ invarie $f^E$

\margin{\unsure{itou|i.e.}}
% slanted
c) $U_\rho$ \rule{1.6cm}{0.35pt} $Q_X$

Si \deleted{de plus} $\rho=0$, ces conditions équivalent aussi à :

d) $\exists$ \inserted{loc.} section de $U_\rho$ partout $\ne1$ qui invarie $f^E$ (ou encore, $Q_X$)
% archive sheet number: 41
